\section{Introduction: water, soil and agriculture}
Agriculture is the largest user of fresh water: about seventy per cent of the water withdrawn from rivers, lakes and aquifers goes to irrigation, and most of it does not come back to the stream it left. The water that a crop uses passes through one thin layer of the earth, the soil, and the behaviour of that layer decides how much of an applied millimetre is held for the roots, how much is lost to the air and how much drains out of reach. A designer of irrigation systems, a modeller of the soil--plant--atmosphere system and a farmer who reads a tensiometer all need the same small set of ideas: how much water the soil holds, how tightly, and how fast it moves.

This chapter gives those ideas in the order in which a textbook of soil physics gives them: first the soil as a mixture of solids, water and air; then its texture and the classes into which texture divides it; then the amount of water it holds and the ways of measuring that amount; then the energy with which the water is held, which decides where it moves; and last the characteristic constants (saturation, field capacity, wilting point) that turn all this into numbers for a schedule. The second half of the chapter, which carries over from the studies on which this book is based, derives the retention and conductivity functions that the whole of Daisy and of the SWAFL application use. Every definition is followed by a solved example; the chapter ends with a design project, exercises and questions.

\section{The soil as a three-phase system}
A soil is a porous body made of three phases: solid particles (mineral grains and organic matter), liquid water with dissolved salts, and gas (air with water vapour). If a sample of total volume $V_t$ is divided into the volume of solids $V_s$, of water $V_w$ and of air $V_a$, the volume of the pores is $V_v=V_w+V_a$ and $V_t=V_s+V_v$. The masses are $M_s$ (dry solids) and $M_w$ (water); the mass of air is neglected. The following quantities are defined from these six numbers.

\begin{equation}\rho_b=\frac{M_s}{V_t},\qquad \rho_s=\frac{M_s}{V_s},\qquad \phi=\frac{V_v}{V_t}=1-\frac{\rho_b}{\rho_s},\qquad e=\frac{V_v}{V_s}=\frac{\phi}{1-\phi}\label{eq:phase}\end{equation}


Here $\rho_b$ is the dry bulk density, $\rho_s$ the particle density (about 2.65 Mg\,m$^{-3}$ for mineral soils), $\phi$ the porosity and $e$ the void ratio. The water content is described in two ways. The gravimetric content $w=M_w/M_s$ is a mass ratio and is what a laboratory oven measures; the volumetric content $\theta=V_w/V_t$ is what a model uses, because it multiplies a depth. The two are related by $\theta=w\,\rho_b/\rho_w$, with $\rho_w=1$ Mg\,m$^{-3}$. The degree of saturation is $S=V_w/V_v=\theta/\phi$ and the air-filled porosity is $\phi_a=\phi-\theta$.

\exbox{Phase relations of a core}{An undisturbed core of volume $V_t=250$ cm$^3$ has a wet mass of 402 g and, after drying at 105\,$^\circ$C, a dry mass of 337.5 g. Take $\rho_s=2.65$ Mg\,m$^{-3}$. Find the dry bulk density, the porosity, the void ratio, the gravimetric and volumetric water contents, the degree of saturation and the air-filled porosity.}\sol Given: $V_t=250$ cm$^3$, $M_{wet}=402$ g, $M_s=337.5$ g. \subhd{Bulk density.} $\rho_b=337.5/250=1.35$ g\,cm$^{-3}$ \ans \subhd{Porosity and void ratio.} $\phi=1-1.35/2.65=0.4906$ and $e=0.4906/(1-0.4906)=0.963$ \ans \subhd{Water contents.} $M_w=402-337.5=64.5$ g, so $w=64.5/337.5=0.1911$. Then $\theta=w\rho_b=0.1911\times1.35=0.258$ cm$^3$\,cm$^{-3}$ \ans \subhd{Saturation and air.} $S=0.258/0.4906=0.526$ and $\phi_a=0.4906-0.258=0.233$ \ans So about half of the pore space holds water and a quarter of the soil volume is air.

\section{Soil texture and the classes of soils}
The mineral part of a soil is a mixture of particles of different sizes. By the USDA convention, \kw{sand} is the fraction between 2 and 0.05 mm, \kw{silt} between 0.05 and 0.002 mm and \kw{clay} below 0.002 mm; the proportions of the three (summing to 100 per cent) are the \kw{texture}. The texture triangle divides all possible mixtures into twelve classes. In irrigation the classes are often grouped in three types: \emph{coarse} soils (sand, loamy sand, sandy loam), which take in water fast, hold little and drain quickly; \emph{medium} soils (loam, silt loam, silt), which hold the most plant-available water; and \emph{fine} soils (clay loam, silty clay loam, silty clay, clay), which take in water slowly, hold it tightly and are prone to runoff and waterlogging.

The ten soils used throughout this book are the ones of the Daisy ensemble: S01 coarse loam, S02 coarse sand, S03 sand, S04 loamy sand, S05 sandy loam, S06 loam, S07 silt loam, S08 sandy clay loam, S09 clay loam and S10 silty clay loam. Table~\ref{tab:ten} lists their composition and van Genuchten parameters; the parameter $\alpha$ is, roughly, the inverse of the air-entry suction and $n$ describes the width of the pore-size distribution (a large $n$ is a uniform pore size, as in sand).

\begin{table}[!htbp]\centering\small\caption{The ten soils of the ensemble: composition (per cent by mass), van Genuchten parameters and plant-available water per centimetre of depth.}
\begin{tabular}{llrrrrrr}\toprule
\hd{Soil} & \hd{Class} & \hd{Sand \%} & \hd{Silt \%} & \hd{Clay \%} & \hd{$\alpha$ (cm$^{-1}$)} & \hd{$n$} & \hd{AWC (mm cm$^{-1}$)}\\\midrule
S01 & Coarse loam & 45 & 35 & 20 & 0.0961 & 1.353 & 1.5\\
S02 & Coarse sand & 88 & 8 & 4 & 0.0536 & 1.600 & 0.5\\
S03 & Sand & 92 & 5 & 3 & 0.1450 & 2.680 & 0.6\\
S04 & Loamy sand & 82 & 12 & 6 & 0.1240 & 2.280 & 0.9\\
S05 & Sandy loam & 65 & 25 & 10 & 0.0750 & 1.890 & 1.2\\
S06 & Loam & 40 & 40 & 20 & 0.0360 & 1.560 & 1.7\\
S07 & Silt loam & 20 & 65 & 15 & 0.0200 & 1.410 & 2.0\\
S08 & Sandy clay loam & 55 & 15 & 30 & 0.0590 & 1.480 & 1.4\\
S09 & Clay loam & 30 & 35 & 35 & 0.0190 & 1.310 & 1.6\\
S10 & Silty clay loam & 10 & 55 & 35 & 0.0100 & 1.230 & 1.7\\\bottomrule\end{tabular}\end{table}

\section{The water content of a soil and how it is measured}
The amount of water in a soil can be measured directly or inferred. The methods fall into four types. \kw{Gravimetric sampling} weighs a sample before and after oven drying; it is the reference method but destroys the sample and gives one point in time. \kw{Neutron scattering} counts the slow neutrons that return to a probe after hydrogen nuclei in the soil have thermalised them; it integrates a sphere of 15--25 cm radius. \kw{Dielectric methods} (time-domain reflectometry, capacitance and frequency-domain sensors) measure the dielectric permittivity of the soil, which is about 80 for water and 3--5 for dry soil, and convert it to $\theta$ by the Topp relation $\theta=-5.3\times10^{-2}+2.92\times10^{-2}\varepsilon-5.5\times10^{-4}\varepsilon^2+4.3\times10^{-6}\varepsilon^3$. \kw{Remote sensing} uses the microwave and thermal signal of the surface, and reaches only the top few centimetres.

A water content is turned into a \emph{depth} by multiplying by the thickness: a layer of thickness $z$ holds the depth $d=\theta z$ of water, in the same unit as $z$. Depths of water are the language of irrigation, because rain, evaporation and irrigation are all written in millimetres; one millimetre over one hectare is 10 m$^3$.

\exbox{Water in a profile}{A profile is sampled in four layers. Gravimetric contents on the sampling day: 0--20 cm, $w=0.140$ ($\rho_b=1.30$); 20--40 cm, $w=0.165$ ($\rho_b=1.38$); 40--80 cm, $w=0.190$ ($\rho_b=1.45$); 80--100 cm, $w=0.205$ ($\rho_b=1.50$ g\,cm$^{-3}$). Find the volumetric content of each layer, the depth of water held in each, and in the whole profile.}\sol For each layer $\theta=w\rho_b$ and $d=\theta z$. layer 0--20: $\theta=0.140\times1.30=0.182$, $d=0.182\times200=36.4$ mm; layer 20--40: $\theta=0.165\times1.38=0.228$, $d=0.228\times200=45.5$ mm; layer 40--80: $\theta=0.190\times1.45=0.275$, $d=0.275\times400=110.2$ mm; layer 80--100: $\theta=0.205\times1.50=0.307$, $d=0.307\times200=61.5$ mm. The sum is $d=253.6$ mm \ans

\section{Soil water potential}
Water in a soil has energy, and it moves from where its energy is high to where it is low. The \kw{total potential} $\Psi_t$ is the work needed to move a unit quantity of pure free water from a reference pool to the point in the soil; per unit \emph{weight} it is a length, the \kw{hydraulic head} $H$. Four components are distinguished: the \emph{matric} (capillary and adsorptive) potential $h$, always negative in an unsaturated soil; the \emph{gravitational} potential $z$, the height above a datum; the \emph{osmotic} potential, due to dissolved salts, which acts on root uptake but not on flow in the soil; and the \emph{pressure} potential, positive below a water table. Without salt, $H=h+z$ and flow is driven by the gradient of $H$ (Chapter~2). It is convenient to speak of the \kw{suction} $|h|$ (in cm of water, in kPa with 1 kPa $\approx 10.2$ cm, or as $\mathrm{pF}=\log_{10}|h|$).

\exbox{Heads and flow direction}{A tensiometer at 30 cm depth reads a suction of 250 cm of water, and one at 60 cm reads 210 cm (a datum at the surface, $z$ positive upward). Find the total head at each depth and the direction and size of the flow if the mean conductivity between them is $K=0.8$ cm\,d$^{-1}$.}\sol Heads: $H_{30}=h+z=-250+(-30)=-280$ cm and $H_{60}=-210+(-60)=-270$ cm. \subhd{Direction.} Water moves from high head to low head, that is from 60 cm ($-270$) to 30 cm ($-280$): it moves \emph{upward}. \subhd{Flux.} Darcy's law gives $q=-K\,\Delta H/\Delta z=-0.8\times[(-280)-(-270)]/(30-60)$ with $z$ taken downward positive: the gradient is $10/30=0.333$, so $q=0.8\times0.333=0.267$ cm\,d$^{-1}$ upward \ans This is 2.7 mm per day, comparable to the evaporative demand: the water table is not needed to explain it.

\section{Characteristic soil water constants}
Because $\theta$ and $h$ are linked by the retention curve (later in this chapter), a few suctions have been given names. \kw{Saturation} ($h=0$) is the content at which all pores are filled. \kw{Field capacity} is the content that remains when drainage has become negligible, two or three days after a soaking, and is placed at $h=100$ cm (pF 2.0) in this book (some authors use 330 cm for fine soils). The \kw{permanent wilting point} is the content at which a crop cannot recover its turgor, placed at 15\,849 cm (pF 4.2, 1.55 MPa). The difference is the \kw{plant-available water}, and over the depth $Z_r$ of the roots the \kw{total available water} is

\begin{equation}\mathrm{TAW}=1000\,(\theta_{FC}-\theta_{WP})\,Z_r\qquad(\text{mm, with }Z_r\text{ in m}).\label{eq:taw}\end{equation}
\par\noindent where ${\rm TAW}$ is the total available water (mm), $\theta_{FC}$ and $\theta_{WP}$ the water contents at field capacity and wilting point (m$^3$\,m$^{-3}$), $Z_r$ the rooting depth (m) and 1000 converts m to mm.



A crop is allowed to use a fraction $p$ of the TAW without stress; the \kw{readily available water} is $\mathrm{RAW}=p\,\mathrm{TAW}$, and the depth that has been used is the \kw{depletion} $D_r$. Irrigation is called when $D_r$ reaches RAW and refills it to field capacity, so the net irrigation depth is $Z_n=D_r$ and the interval is $Z_n/\mathrm{ET}_c$.

\exbox{Soil water constants of the loam}{For the loam S06 ($\alpha=0.036$ cm$^{-1}$, $n=1.56$, $\theta_s=0.468$, $\theta_r=0.02$) find the water contents at field capacity (100 cm) and wilting point (15\,849 cm), the plant-available water per metre of soil, and the TAW and RAW ($p=0.5$) for a root depth of 0.8 m.}\sol With $m=1-1/n=0.359$. \subhd{Field capacity.} $(\alpha h)^n=(3.6)^{1.56}=7.38$, so $S_e=(1+7.38)^{-0.359}=0.466$ and $\theta_{FC}=0.02+0.448\times0.466=0.229$. \subhd{Wilting point.} $(\alpha h)^n=(570.6)^{1.56}=19945$ so $S_e=0.0286$ and $\theta_{WP}=0.033$. \subhd{Available water.} $\theta_{FC}-\theta_{WP}=0.196$, that is 196 mm per metre \ans \subhd{TAW and RAW.} $\mathrm{TAW}=1000\times0.196\times0.8=157$ mm and $\mathrm{RAW}=0.5\times157=78$ mm \ans

